% Native-ready author proof.  This file does not modify an inherited proof.
% Dependencies: 260, 281--284, 317, 319, 320; the class is the enlarged class of 284.
\section{Complete canonical collision-boundary budgets}
\label{sec:cs-canonical}

We now control the actual coordinate sides of the original section at
all collision indices.  The source is a genuine collision state, and
all propagation uses its actual past and future flights.  In particular,
the quantity $\tau+R\cos\varphi$ below is an optical coefficient;
it is not an extra segment of the physical roof.

Write $t_*=\tau_-=3/50$, $k_-=100/47$, $k_+=20/9$,
$U=170/9$, and $\sigma=47/59$.  Throughout this section retain the
endpoint cap $|p_0|,|p_m|\le9/10$ whenever a Gauss-angle charge is
estimated.  Put $c_\circ=\sqrt{19}/10$.  The source and root-supremum
estimates themselves do not require this endpoint cap.  Use precisely
the fixed enlarged class
\[
 \widehat B_{\rm cv}\ge\max\{B_{\rm cv},K_{\rm mom},2C_*\},
 \qquad \widehat k_0,\quad \widehat\ell_*,\quad\widehat C_g
\]
of \eqref{eq:mc-class}.  These are its own growth and subdivision
constants; a constant for the smaller marked-foot class is not reused.
Keep $J_f,J_{\rm pre},J_{\rm len},L_0,A_{\rm mom}$ from
\eqref{eq:uc-candidates} and \eqref{eq:mc-pre}--\eqref{eq:mc-A}.
Define
\begin{align}
 A_{\rm fan}
   &=\frac{J_fJ_{\rm len}}{t_*}
                    \left(1+\frac2{\widehat k_0-1}\right),&
 E_{\rm cs}&=\frac1{k_-t_*}=\frac{47}{6},
                                                        \label{eq:cs-source-constants}\\
 H_{\rm hor}
   &=\max\{J_{\rm pre}J_f E_{\rm cs},
              A_{\rm mom}\widehat C_gJ_f E_{\rm cs}\},&
 H_{\rm ver}
   &=\max\{J_fE_{\rm cs},
              A_{\rm fan}\widehat C_gJ_f E_{\rm cs}\}.\notag
\end{align}
The symbols ``horizontal'' and ``vertical'' refer to the unchanged
$(\alpha,p)$ collision coordinates: $p=b$ and $\alpha=a$,
respectively.  The vertical source is parametrized by
$\varphi=\arcsin p$, without changing its physical states.

\subsection{Actual first images and one-sided supremum sums}

\begin{lemma}[The two canonical sources in the qualified curve class]
\label{lem:cs-first}
Fix a collision disk and translate its center to the origin.
For $b\in(-1,1)$ the horizontal source $p=b$, parametrized by
$\alpha$, has the complete first-image family of
Lemma~\ref{lem:mc-family}.  For every fixed $a$ the vertical source
$\alpha=a$, parametrized by $\varphi\in(-\pi/2,\pi/2)$,
has a complete countable family of short homogeneous first images in
$\mathcal W_R^u(\widehat B_{\rm cv})$.  Its exact pushed densities
satisfy
\begin{equation}\label{eq:cs-fan-density-sum}
                    \sum_C\|\rho_C\|_\infty\le A_{\rm fan}.
\end{equation}
Both families are source-disjoint and cover their source up to its
actual singular points and a countable set of subdivision points.
Their total source lengths are $2\pi$ and $\pi$, respectively;
neither source is normalized to a probability.
\end{lemma}
\begin{proof}
The horizontal statement, including its real first-hit test and
uniformity as $b$ tends to either grazing value, is exactly
Lemmas~\ref{lem:mc-first} and \ref{lem:mc-family}.  Here is also
a common first-image calculation that fixes the optical conventions
needed below.  Put $c=\sqrt{1-b^2}$ on the horizontal source.
Its outgoing direction is $\vartheta=\alpha+\arcsin b$, and
its current contact is
\[
 q_0=-Rb\,v_\vartheta^\perp+Rc\,v_\vartheta.
\]
Consequently the next actual contact can be written
$q_1=-Rb\,v_\vartheta^\perp+a_1v_\vartheta$, where
$a_1=\tau+Rc\ge t_*$.  The line foot lies inside the current
disk in general.  The flight starts at the actual exit $q_0$;
one must not apply a first-hit test starting from that foot.
For a vertical source $q_0$ is fixed, and
$q_1=q_0+a_1v_\vartheta$ with $a_1=\tau\ge t_*$.
In the following identities take $b_1=Rb$ in the horizontal case
and $b_1=0$ in the vertical case, and let $d=\cos\varphi_1$:
\begin{align}
 \left|\frac{ds_1}{d\vartheta}\right|&=a_1/d,&
 V_1&=\kappa+d/a_1,&
 \rho_1&=\frac{d}{a_1\sqrt{1+V_1^2}},
                                                        \label{eq:cs-first-identities}\\
 (a_1)_{s_1}&=p_1+b_1d/a_1,&
 (V_1)_{s_1}
   &=-\frac{p_1\kappa}{a_1}
        -\frac{2p_1d}{a_1^2}-\frac{b_1d^2}{a_1^3}.\notag
\end{align}
The horizontal identities agree with \eqref{eq:mc-first}--
\eqref{eq:mc-density}, since $a_1=R(c+\tau\kappa)$.
For the vertical source they follow by differentiating
$q_1=q_0+\tau v_\vartheta$ and projecting onto the target tangent
and the flight direction.  Thus both have $k_-\le V_1\le U$ and
\[
 |(V_1)_{s_1}|
  \le\frac{k_++2/t_*}{t_*}+\frac{R_+}{t_*^3}
  =K_{\rm mom}.
\]
The vertical bound is smaller, since its last term is zero.  This
verifies membership in the already enlarged class before invoking
any growth estimate.

For completeness the vertical source also has a uniform finite
physical first-hit partition.  Its current point is exterior to every
other disk, with $|z-q_0|\ge1-R_+>R_+$.  The directions that enter
one other disk form its open angular shadow of width less than
$\pi$.  Intersecting it with the outgoing semicircle gives an interval,
with at most two tangency endpoints.  The actual first contact has
$0<|z|\le D_f$, so at most $2N_f$ such endpoints must be cut.
Distinct entry times cannot exchange order without a common point of
two disjoint disks.  Hence there are at most $J_f=2N_f+1$ first-hit
intervals.  Outward directions cannot reenter the current convex disk.
Every flight being partitioned has the true lower bound $\tau\ge t_*$.

On each first image $V_1>0$, so $\varphi_1$ is strictly monotone and
each homogeneous strip is met at most once.  Its total length is at
most $L_0$.  Apply the fixed length subdivision with at most
$J_{\rm len}$ pieces per strip.  In the central strip
$\rho_1\le t_*^{-1}$; in each signed strip of index $k$ it is at
most $(t_*k^2)^{-1}$.  Summing the actual densities, rather than
counting infinitely many strips without their weights, proves
\eqref{eq:cs-fan-density-sum}.  All collision angles on a connected image are represented by a
continuous unwrapped lift, so crossing the numerical value $0$ creates
no new physical interval.  Such a lift exists on its source interval;
the monotone target $\varphi_1$ and the bound $L_0$ concern that
same lift.  The horizontal source itself is opened at one common fixed
seam before its preliminary arcs are chosen; the retained $J_{\rm pre}$
includes that division.  The vertical source is already an interval.
Thus $J_f$ has only its proved physical tangency cuts.  If a short-curve
implementation additionally splits an image at a coordinate seam,
the one possible split is included in
$J_{\rm len}=2+\lceil2L_0/\widehat\ell_*\rceil$, exactly the
constant of \eqref{eq:mc-length}, not silently charged to $J_f$.
There is at most one such seam crossing because the lifted
horizontal span is at most $\pi/k_-<2\pi R_-$; after the fixed
length subdivision, adding one cut fits within that displayed
conservative count.  In the one-flight case and at the
uncut terminal image we use the continuous lift and do not insert
that extra split.  All removed tangencies
are genuine singular endpoints, and the remaining subdivision cuts
are countable.  They have zero source length.  No true grazing
collision is deleted from an admissible physical record.
\end{proof}

For $l\ge1$, propagate a source variation across exactly $l$ actual
flights, including the $l-1$ intermediate reflection kicks and neither
an initial nor a terminal kick.  In unfolded transverse coordinates
let $M_l$ be this positive optical matrix and put
\begin{equation}\label{eq:cs-rows}
 (X_l,Y_l)=(\kappa_l,c_l)M_l,\qquad
 A_l^{\rm hor}=RcX_l+Y_l,\qquad A_l^{\rm ver}=Y_l.
\end{equation}
These are the absolute derivatives of the actual terminal momentum
with respect to $\alpha$ on a horizontal source, or $\varphi$ on
a vertical source.  A common orientation sign is suppressed, not an
absolute value inside either positive sum.

\begin{theorem}[Complete branch-supremum bounds for canonical sources]
\label{thm:cs-half-sup}
For every $R$, $l\ge1$, and either canonical source there is a
countable partition into disjoint open source intervals $\mathcal J_l$.
It retains every actual regular collision history, up to the null
cuts just described.  The first $l-1$ images are homogeneous and
short according to the growth partition; the terminal image is not
artificially divided into infinitely many strips.  On these intervals,
\begin{align}
 \sum_{I\in\mathcal J_l^{\rm hor}}
                 \sup_I(A_l^{\rm hor})^{-1}&\le H_{\rm hor},&
 \sum_{I\in\mathcal J_l^{\rm ver}}
                 \sup_I(A_l^{\rm ver})^{-1}&\le H_{\rm ver},
                                                        \label{eq:cs-half-sums}\\
 \sup_I(A_l^{\rm hor})^{-1},\quad
 \sup_I(A_l^{\rm ver})^{-1}
                  &\le E_{\rm cs}\sigma^{l-1}.\notag
\end{align}
These are suprema on whole intervals, not a root sum at a common
target momentum.  They hold without a terminal incidence cap and
in the reversed physical direction with the same constants.
\end{theorem}
\begin{proof}
The initial transverse variation is $(Rc,1)$ for a horizontal
source and $(0,1)$ for a vertical source.  These conventions also
follow directly from the physical differential of
Lemma~\ref{lem:eo-differential}.  If
\[
 B(c)=\begin{pmatrix}c&0\\\kappa&1\end{pmatrix},\quad
 P(\tau)=\begin{pmatrix}1&\tau\\0&1\end{pmatrix},\quad
 K(c')=\begin{pmatrix}1&0\\2\kappa/c'&1\end{pmatrix},
\]
then its positive one-flight matrix $\mathsf D$ satisfies
$B(c')\mathsf D=K(c')P(\tau)B(c)$.  At the terminal point
$c_l(0,1)B(c_l)^{-1}K(c_l)=(\kappa_l,c_l)$.
Thus multiplication leaves exactly the $l$ real flights and
$l-1$ intermediate kicks in $M_l$, with no terminal kick left
over.  The initial canonical vectors $(R,0)$ and $(0,1)$ are
sent by $B(c)$ to $(Rc,1)$ and $(0,1)$, respectively.
After the first true flight
its positive displacement is $B_1=Rc+\tau$ or $B_1=\tau$, and
its angular component is $D_1=1$.  Each subsequent actual
intermediate kick and flight acts by
\[
 B_{r+1}=(1+\tau_r k_r)B_r+\tau_r D_r,\qquad
 D_{r+1}=k_rB_r+D_r,\qquad
 k_r=\frac2{R c_r}.
\]
Therefore $B_l\ge t_*(59/47)^{l-1}$.  The terminal incoming
momentum row is $\kappa_l B_l+c_lD_l\ge k_-B_l$.
This proves the pointwise and hence the branch-supremum optical
bounds, with their stated index $l-1$.  It introduces no virtual
physical flight.

For $l=1$, use just the finite first-hit partitions from
Lemma~\ref{lem:cs-first}.  Their sizes are at most
$J_{\rm pre}J_f$ and $J_f$, giving the first alternatives in
\eqref{eq:cs-source-constants}.  For $l\ge2$, push the exact
first-image densities $\rho_C$ through $l-2$ actual iterates.
On every resulting short homogeneous graph $V$, the final true
collision has at most $J_f$ physical branches and obeys the exact
last-collision cancellation in the proof of
Theorem~\ref{thm:mc-kernel}:
\[
       \left|\frac{dp_l}{d\ell_{l-1}}\right|\ge k_-t_*.
\]
For a source interval leading to one such final branch the chain
rule therefore bounds the inverse source derivative by
\[
 \|\rho_C\|_\infty\,
          \sup_VJ_V(T_R^{-(l-2)})\,E_{\rm cs}.
\]
The growth theorem for the qualified enlarged class bounds the sum
of the middle factor by $\widehat C_g$.  Sum first over the at most
$J_f$ terminal physical branches and then over the full weighted
first-image family.  This proves the second alternatives in
\eqref{eq:cs-source-constants}.  No supremum is transferred from a
pointwise endpoint kernel to a collection of branches.

The physical and homogeneous partitions and their monotone pullbacks
are countable and source-disjoint.  The terminal momentum is strictly
monotone on each one.  A terminal interval restriction therefore
leaves at most one interval and does not increase any of the sums.
A common fixed source seam is used before superposing two families.
Time reversal $\mathcal I(\alpha,\varphi)=(\alpha,-\varphi)$
is an isometry on both source parameters; it preserves every bound
and transports the actual source tests.  The physical map is
invertible on each regular history.  Translating the selected
collision disk to the origin consequently creates no multiplicity
of initial points, lattice labels, or collision words.
\end{proof}

\subsection{Exact two-sided determinants, including the incidence factor}

\begin{lemma}[Canonical two-sided root weights]
\label{lem:cs-determinant}
At an actual interior collision $i\in\{1,\ldots,m-1\}$ write
$c_i=\sqrt{1-p_i^2}$ and let $(X_-,Y_-)$ and $(X_+,Y_+)$ be
\eqref{eq:cs-rows} for the $i$ backward and $m-i$ forward flights.
Put $A_\pm^{\rm hor}=Rc_iX_\pm+Y_\pm$,
$A_\pm^{\rm ver}=Y_\pm$, and
\[
 d_\pm=\frac{Y_\pm}{Rc_iX_\pm+Y_\pm},\qquad
 h_\pm=1+\frac{Rc_iX_\pm}{Y_\pm},\qquad
 C_0=\frac{t_*+R_+}{2t_*}=\frac{53}{12}.
\]
The unchanged endpoint root weight
$u_x=|\partial_{\alpha_0}p_m|^{-1}$ satisfies
\begin{align}
 u_x&=\frac{c_i}{A_-^{\rm hor}A_+^{\rm hor}(d_-+d_+)}
       \le\frac{C_0c_i}{A_-^{\rm hor}A_+^{\rm hor}},
                                                        \label{eq:cs-horizontal-u}\\
 u_x&=\frac{c_i}{A_-^{\rm ver}A_+^{\rm ver}(h_-+h_+)}
       \le\frac{c_i}{2A_-^{\rm ver}A_+^{\rm ver}}
       \le\frac1{2A_-^{\rm ver}A_+^{\rm ver}}.
                                                        \label{eq:cs-vertical-u}
\end{align}
In particular the horizontal bound retains the factor
$c_i=\sqrt{1-b^2}$ uniformly as its level approaches grazing.
\end{lemma}
\begin{proof}
Up to a common sign in each endpoint row, the derivative columns
in the central variables $(\alpha_i,\varphi_i)$ are
\[
 \begin{pmatrix}
       A_-^{\rm hor}&-Y_-\\
       A_+^{\rm hor}& Y_+
 \end{pmatrix}.
\]
Indeed varying $\alpha_i$ with $\varphi_i$ fixed gives the
initial transverse field $(Rc_i,1)$ in both physical directions.
Varying $\varphi_i$ at fixed contact gives $(0,1)$ forward and
$(0,-1)$ backward.  Time reversal may reverse the entire backward
momentum row but cannot change this relative sign between columns.
The determinant thus has absolute value
$A_-^{\rm hor}Y_++Y_-A_+^{\rm hor}$.

Collision-area invariance and the exact endpoint coarea identity give
\[
 d\nu=\frac{d\alpha_i\,dp_i}{4\pi}
      =\frac{c_i\,d\alpha_i\,d\varphi_i}{4\pi}
      =\frac{u_x}{4\pi}\,dp_0\,dp_m.
\]
This proves both displayed identities; unlike a marked line foot,
a canonical collision source has no additional factor $R^{-1}$.
For the inequalities, write $M_l=\widetilde M P(\tau_1)$, with
$P(\tau)=\left(\begin{smallmatrix}1&\tau\\0&1\end{smallmatrix}\right)$.
If $(x,y)=(\kappa_l,c_l)\widetilde M$, positivity implies
$X_l=x$, $Y_l=\tau_1x+y$, and hence
\[
       0< X_l/Y_l\le1/t_*.
\]
Therefore $d_\pm\ge t_* /(t_*+R_+)$ and $h_\pm\ge1$.
The coefficient $c_i$ arises from the original collision area and
has not been replaced by one in \eqref{eq:cs-horizontal-u}.
\end{proof}

\begin{corollary}[All-word two-sided supremum budgets]
\label{cor:cs-overlap}
Intersect the two capped or uncapped one-sided partitions at the
same canonical source.  Let $\mathcal O_i^{\rm hor}$ or
$\mathcal O_i^{\rm ver}$ denote the nonempty interval overlaps.
Then
\begin{align}
 \sum_{O\in\mathcal O_i^{\rm hor}}\sup_O u_x
   &\le C_0c_i H_{\rm hor}E_{\rm cs}
                    (\sigma^{i-1}+\sigma^{m-i-1}),
                                                        \label{eq:cs-overlap-h}\\
 \sum_{O\in\mathcal O_i^{\rm ver}}\sup_O u_x
   &\le\tfrac12H_{\rm ver}E_{\rm cs}
                    (\sigma^{i-1}+\sigma^{m-i-1}).
                                                        \label{eq:cs-overlap-v}
\end{align}
The first assertion is uniform over all $b\in(-1,1)$, and the
second over all fixed contact angles.  No exact-record selector is
used to form new interval components.
\end{corollary}
\begin{proof}
Apply the weighted countable interval-overlay lemma
\ref{lem:mb-overlay} to the two actual source-disjoint families.
For clarity its bound is
$\sum_{I\cap J\ne\varnothing}a_Ib_J
 \le(\sup_Jb_J)\sum_Ia_I+(\sup_Ia_I)\sum_Jb_J$.
It follows by assigning an overlap to the interval whose left endpoint
is later: the other disjoint family has at most one interval containing
that endpoint.  Equal endpoints are assigned once by a fixed rule;
open intervals and finite exhaustion handle all countable cases.
Use \eqref{eq:cs-half-sums} for the full sums and the optical bounds
with half lengths $i$ and $m-i$ for their maxima, then apply
Lemma~\ref{lem:cs-determinant}.  On a horizontal level $c_i$ is a
constant.  On a vertical side it is at most one.  The partitions
cover all actual regular histories; their artificial endpoints have
zero angular-image length and are not new physical records.
\end{proof}

\subsection{Uniform total turning on the actual source intervals}

Here are explicit variation constants for the same enlarged class.
In the constants of \eqref{eq:ab-constants-one}--
\eqref{eq:ab-constants-three}, replace
$B_{\rm cv},k_0,\ell_*,D_B$ by
$\widehat B_{\rm cv},\widehat k_0,\widehat\ell_*,\widehat D_B$,
where $\widehat D_B$ is \eqref{eq:uc-DB} for that allowance.
Equivalently, put
\begin{align}
 \widehat B_2&=\widehat B_{\rm cv}(1+U^2)^{3/2},\notag\\
 \Gamma&=\sqrt{1+k_+^2},\qquad
 L_h=3\sqrt{1+U^2}/k_-,\notag\\
 \widehat c_H&=\sin(\widehat k_0^{-2}),\notag\\
 \widehat V_H&=\max\{-2\log\widehat c_H,
        \log(\pi/2)+2\log(1+\widehat k_0^{-1})\},\notag\\
 \widehat D_1&=\max\{(U/\widehat c_H+\widehat D_B)L_0^{2/3},
                   (2\pi U+\widehat D_B)L_h^{2/3}\},\notag\\
 \widehat V_J&=\frac{\widehat D_1\Gamma^{1/3}L_0^{1/3}}
                         {1-\sigma^{1/3}},\notag\\
 V_{\rm cap}&=-2\log c_\circ,\notag\\
 V_{\rho,\rm cs}&=\widehat V_H+
       (t_*^{-1}+R_+t_*^{-2}+K_{\rm mom}/2)\widehat\ell_*,\notag\\
 V_{t,\rm cs}&=V_{\rm cap}+\widehat B_2L_0/k_-,\notag\\
 V_{1,\rm cs}&=\frac{L_0(k_+R_++1)}{k_-t_*^2},\notag\\
 V_{A,\rm cs}&=\max\{V_{1,\rm cs},
                  V_{\rho,\rm cs}+\widehat V_J+V_{t,\rm cs}\}.
                                                        \label{eq:cs-variation-constants}
\end{align}
These finite constants precede every radius, collision count,
source level, exact label, and roof supremum.

\begin{lemma}[Logarithmic variation and regular angular turning]
\label{lem:cs-turning}
On every one-sided canonical source interval, after restricting its
terminal momentum to $[-9/10,9/10]$,
\[
       \operatorname{Var}\log A_l^{\rm hor}\le V_{A,\rm cs},
       \qquad
       \operatorname{Var}\log A_l^{\rm ver}\le V_{A,\rm cs}.
\]
On every interior two-sided overlap, with $P=(-p_0,p_m)$, one has
\begin{equation}\label{eq:cs-turning}
             \int_{\{P\ne0\}}|d\arg P|
                             \le 2\pi+V_{A,\rm cs}.
\end{equation}
On a fixed initial-angle side, including both endpoint caps, the
corresponding bound is
\begin{equation}\label{eq:cs-end-turning}
          \int_{\{P\ne0\}}|d\arg P|
                 \le V_{\rm end}:=
                       2\pi+\tfrac12(V_{A,\rm cs}+V_{\rm cap}).
\end{equation}
No critical incoming fan at $P=0$ is assigned zero mass by these
regular, punctured-angle statements.
\end{lemma}
\begin{proof}
All the following estimates concern absolute integrals of derivatives,
not merely differences of their endpoint values.  The proof of
Lemma~\ref{lem:ab-jacobian-variation} applies to the enlarged class:
on a homogeneous image of length $r$ the absolute one-step log-Jacobian
variation is at most $\widehat D_1 r^{1/3}$, and the adapted expansion
gives $L_j\le\Gamma\sigma^{l-j}L_l$.  Summation therefore bounds
the full log-Jacobian variation by $\widehat V_J$.

For $l\ge2$ the actual first-image density in
\eqref{eq:cs-first-identities} gives exactly
\[
 A_l^{\rm hor}\ \hbox{or}\ A_l^{\rm ver}
   =\frac{c_lV_l}{\sqrt{1+V_l^2}}
                           \frac{J_C(T_R^{l-1})}{\rho_C}.
\]
The terminal cap is one interval because $p_l$ is strictly monotone,
and it lies in the central homogeneous strip.  On the first homogeneous
piece the variation of $\log c_1$ is at most $\widehat V_H$.
The exact identities \eqref{eq:cs-first-identities} give
$|(a_1)_{s_1}|\le1+R_+/t_*$ and $|(V_1)_{s_1}|\le K_{\rm mom}$.
Their integrals, with $a_1\ge t_*$, bound
$\operatorname{Var}\log\rho_C$ by $V_{\rho,\rm cs}$.
The terminal factor has variation at most $V_{t,\rm cs}$, since
$\varphi_l$ is monotone, $c_l\ge c_\circ$, and
$|(V_l)_s|\le\widehat B_2$.  This proves the claimed bound for
$l\ge2$, preserving the exact initial density and all its grazing
factors throughout the summation.

For $l=1$ no initial homogeneous cut is needed.  In either source
$A_1=\kappa a_1+d$, and
\[
       \partial_{s_1}A_1=d(\kappa b_1-p_1)/a_1.
\]
Its log-derivative is bounded by
$(k_+R_++1)/(k_-t_*^2)$.  Integrate over a first-image graph of
length at most $L_0$.  This proves the remaining alternative.
Time reversal preserves these absolute variation estimates.

On a two-sided interior overlap the two components of $P'$ have
fixed signs and nonzero magnitudes $A_-^{\rm hor},A_+^{\rm hor}$,
or $A_-^{\rm ver},A_+^{\rm ver}$.  Thus its tangent stays in one
quadrant and has total turning at most
\[
 \tfrac12\{\operatorname{Var}\log A_-+
                    \operatorname{Var}\log A_+\}\le V_{A,\rm cs}.
\]
One component of $P$ is strictly monotone, so $P$ crosses the origin
at most once.  Lemma~\ref{lem:ab-radial-turning} proves
\eqref{eq:cs-turning}, including both punctured sides of a crossing.
At a fixed initial angle use the source variable $\varphi_0$.
Then $P'=(-c_0,\pm Y_m)$, its tangent remains in one quadrant, and
$\operatorname{Var}\log c_0\le V_{\rm cap}$ on the initial cap.
The same lemma proves \eqref{eq:cs-end-turning}.  A terminal-angle
side is covered by time reversal, which preserves absolute turning.
\end{proof}

\subsection{Complete section-side charges, including the endpoint indices}

For one actual regular canonical side let its charge mean the sum,
over its unique physical source points and all words, of
$A|d\theta|$, where
\[
 A=\frac{\omega}{4\pi c_*\chi},\qquad
 \theta=\arg(-p_0,p_m),\qquad
 P=(-p_0,p_m)\ne0.
\]
An additional existing selector of absolute value at most one can
restrict this positive integral.  This definition counts the physical
side once.  When a piecewise record weight is differentiated, its
jump multiplies this one-copy charge; the two incident traces are
combined into that jump before absolute values.  If two incident
traces are intentionally counted separately, twice this bound is
available.  It is not legitimate to count them twice and then omit
the factor two.

Set
\begin{align}
 F_{\rm cs}&=\frac{\Lambda_\circ}{4\pi c_*R_-^2},&
 C_{\rm hor}&=F_{\rm cs}(2\pi+V_{A,\rm cs})
                         C_0H_{\rm hor}E_{\rm cs},\notag\\
 C_{\rm ver}&=\tfrac12F_{\rm cs}(2\pi+V_{A,\rm cs})
                                      H_{\rm ver}E_{\rm cs},&
 E_{\rm hor}&=\pi F_{\rm cs}C_{\rm mom},\notag\\
 E_{\rm ver}&=F_{\rm cs}V_{\rm end}H_{\rm ver}.
                                                        \label{eq:cs-charge-constants}
\end{align}

\begin{theorem}[Uniform canonical side charges on the physical source]
\label{thm:cs-charge}
For every $R\in[R_-,R_+]$, $m\ge2$, and $1\le i\le m-1$,
the complete regular charge at $p_i=b\in(-1,1)$ or at
$\alpha_i=a\pmod{2\pi}$ satisfies respectively
\begin{align}
 \mathfrak S^{\rm hor}_{i,b}
    &\le C_{\rm hor}\sqrt{1-b^2}
                         (\sigma^{i-1}+\sigma^{m-i-1}),
                                                        \label{eq:cs-charge-h}\\
 \mathfrak S^{\rm ver}_{i,a}
    &\le C_{\rm ver}(\sigma^{i-1}+\sigma^{m-i-1}).
                                                        \label{eq:cs-charge-v}
\end{align}
The vertical source may be restricted to the actual momentum interval
of an original rectangle, or to any positive Borel subsource, without
an added component count.  At either endpoint index $0$ or $m$ a
horizontal level in the endpoint cap has charge at most
$E_{\rm hor}$, and a vertical level has charge at most $E_{\rm ver}$.
These endpoint assertions hold also for $m=1$.

In particular each fixed type and level of interior side has total
charge at most $2C/(1-\sigma)$ after all its collision indices are
summed.  The indices satisfying
$\min(i-1,m-i-1)\ge D$ have total charge at most
$2C\sigma^D/(1-\sigma)$.  These are sums over complete physical
histories before any count, radius, label or source-level supremum.
\end{theorem}
\begin{proof}
The exact optical identity of Corollary~\ref{cor:bf-wronskian} gives
\[
 A=\frac{u_x}{4\pi c_*h_x},\qquad
 h_x=R^2e^{\mathsf T}G^{-1}e\ge R_-^2/\Lambda_\circ.
\]
It follows that $A\le F_{\rm cs}u_x$.
On each interior overlap use Lemma~\ref{lem:cs-turning} and then
Corollary~\ref{cor:cs-overlap}.  Nonnegative countable summation
gives \eqref{eq:cs-charge-h}--\eqref{eq:cs-charge-v}.  A Borel
restriction is made after this positive summation; it is not
assumed to have a finite number of components or a differentiable
boundary.  The positive source density, the conditional factor
$c_*^{-1}$ and the original collision-area normalization have all
been retained.

For an initial horizontal endpoint side $p_0=b\ne0$, parametrize
by the actual terminal momentum $v=p_m$.  Its absolute angular
differential is $|b|(b^2+v^2)^{-1}dv$.  The pointwise physical
root sum of Theorem~\ref{thm:mc-kernel}, evaluated at the actual
pair $(b,v)$, gives the complete charge bound
\[
 F_{\rm cs}C_{\rm mom}
       \int_{-9/10}^{9/10}\frac{|b|}{b^2+v^2}\,dv
                   \le\pi F_{\rm cs}C_{\rm mom}=E_{\rm hor}.
\]
For $b=0$ its regular polar angle is constant on either side of
$v=0$, so the regular angular charge is zero.  Its critical incoming
source is a separate endpoint source and is not discarded.
Time reversal proves the terminal horizontal assertion.

For a fixed initial angle, on each forward fan interval
\[
 u_x=\frac1{Rc_0X_m+Y_m}\le\frac1{Y_m}.
\]
The branch-supremum sum in Theorem~\ref{thm:cs-half-sup} is therefore
at most $H_{\rm ver}$.  Multiply by $F_{\rm cs}$ and the complete
turning bound \eqref{eq:cs-end-turning}.  This proves the stated
vertical endpoint bound, and time reversal proves the final one.
The original cap and any rectangle-side momentum interval merely
restrict these positive charges.  A coordinate seam, subdivision
endpoint or section corner has zero angular-image length and creates
no additional charge.  All visited grazing limits and critical
incoming fans retain their separately proved physical traces.

Finally sum the two geometric sequences in the interior bounds.
The deep-index restriction starts each surviving exponent at least
$D$, giving the asserted exponential tail.  For every fixed $R,m$
all local source limits and countable exhaustions precede these
uniform estimates.  No long-time or outer-cutoff limit was interchanged.
\end{proof}

\paragraph{A useful grazing-level consequence.}
The factor in \eqref{eq:cs-horizontal-u} and
\eqref{eq:cs-charge-h} is essential rather than cosmetic.
For the original signed homogeneous levels
$\varphi=\pm(\pi/2-k^{-2})$ it is
$\cos\varphi=\sin(k^{-2})\le k^{-2}$.  Thus their complete
horizontal charges are summable over $k$, even though there are
infinitely many such levels.  This consequence does not assert
that every whole-word cell has already been assigned to such a
boundary; any use requiring that assignment must prove it separately.
